Nobody was in charge, and by the second hour after dawn a considerable number of people had come to Tyrion Lannister to find out who it was.

That was the whole of the first day and he understood it before the sun had properly cleared the Dragonpit. There was no king. There was no queen, or rather there was one, on a pallet in the Great Hall with an arm broken in two places and a great deal of internal damage nobody wanted to describe out loud, who had not been declared anything by anybody and whose army was standing about in the streets waiting to be told what it was for. There was no Small Council, because the Small Council had been Cersei’s and consisted at present of one dead maester and a Master of Whisperers at the bottom of a cellar stair. There was no Lord Commander of the City Watch, because the City Watch had ceased to exist as an institution somewhere around the ninth hour of the second day. There was no High Septon; the last one had died in the Sept’s ruin two years past and the Faith had not agreed on a replacement.

What there was, was Tyrion Lannister, who had been Hand of the King to a boy who had him arrested for murder, and Hand of the Queen to a woman who was currently unconscious, and who possessed no standing of any kind that anyone could have written down.

“My lord,” said the third man that morning. He was a guild steward and he had a list. “The Street of Steel. Who’s to say which forges reopen? Half of them have no masters left living, and there’s men squatting in two of them already saying possession’s nine-tenths of it, and we’ll have blood before noon if somebody doesn’t rule on it.”

“On what authority would I rule on it?”

The steward looked at him with the flat exhausted incomprehension of a man who had walked four miles through a burning city to ask a question and did not want philosophy.

“My lord, there’s nobody else.”

\textit{There it is,} Tyrion thought. \textit{Written on no parchment, sworn by nobody, and entirely true. The whole of government, at the beginning, is one man standing where the question arrives.}

“Forges one through nine to their surviving masters or their masters’ widows,” he said. “Squatters out by dusk, and tell them they’ll get first claim on the tenth through fourteenth, which have no masters at all, once somebody’s established that they can actually smith. If they can’t, they’re carters now, and there’s rather a lot of carting to be done.”

“Yes, my lord. And by whose — “ The steward stopped himself. “Thank you, my lord.”

He went. Tyrion stood in the courtyard and thought: \textit{I have just disposed of nine forges I do not own, on behalf of a state that does not exist, and it will hold, because he will go and say that it was ruled and nobody will ask him by whom.}

He found that observation considerably more frightening than anything that had happened to him during the battle.

\sceneBreak

He went to see about his brother at midmorning, because it was the only item on the day’s business that was his and not the city’s, and because he had been putting it off for eleven hours.

They had Jaime in a storeroom off the kitchens with nine other men, because the kitchens were where the wounded had been and the storeroom was cold. Somebody had done it properly — the cloak folded, the hands arranged, the golden hand set beside him rather than on him, which Tyrion thought about for some time afterward and never worked out whether to be grateful for.

“I’d like him taken to the Great Sept,” he said. “What’s left of it. There’s a Lannister vault at Casterly Rock but there’s no getting him there before spring, so — the Sept, and we’ll move him properly when the roads are open.”

The steward — a different one; there seemed to be an unlimited supply — did not answer immediately, and Tyrion turned and looked at him.

“My lord. The ground’s frozen.”

“Then have it broken.”

“It’s been tried, my lord. On the seventh day of the siege, for the men who died on the works, before — “ He gestured vaguely at the whole of the last two days. “It’s frozen four feet down and it’s been frozen since before the turn of the year. You’d need a fire on the ground for a day and a night to soften enough for one man, and there’s no wood in this city to spare for softening ground.”

“There’s no wood.”

“There’s no wood, my lord. Not a stick. The last of the pier timbers went into the trench-line and everything else went into cookfires in the second month.” The steward’s voice was perfectly level and Tyrion understood he had already had this conversation with a dozen people and had stopped expecting any of them to like it. “There’s no burying and there’s no burning. Not for one man and not for — “

He stopped.

“Say the number,” Tyrion said.

“Nobody has a number, my lord.”

“Say the one you’ve heard.”

“Thirty thousand,” the steward said. “Some say more. In the city and on the works and the ones on the flats, and — “ He swallowed. “And the ones that came up, my lord. Out of the cellars. Those have to go somewhere too and they were already dead once.”

Tyrion stood in a cold storeroom with his brother’s body and nine strangers’ and did the arithmetic he had spent his whole life being good at, and arrived, in about four seconds, at the shape of the thing that was going to define the next month of his life.

\textit{I cannot bury my brother. Nobody in this city can bury anybody. And in nine days, when it warms even slightly, thirty thousand bodies are going to begin, and this city has no clean water and half a granary and a population that has spent two days watching its own dead climb out of the ground.}

“Leave him where he is,” he said. “The cold’s holding. That’s not a solution, it’s a — “ He made himself finish it. “It’s a delay. Go and find Samwell Tarly and Lord Davos and whoever the Faith has put up as its loudest voice this week, and tell them I want them in the Tower of the Hand this afternoon. Not the Great Hall. The Hall’s full of the living.”

\sceneBreak

He walked to the meeting rather than being carried, which took him an hour and a half and was the single most useful thing he did that day.

The Street of Flour had gone. Not damaged — gone, two hundred feet of it, where the reanimated dragon had come down and slid, and what was left was a black glazed trench with the shells of houses standing at either end of it like teeth on both sides of a missing one. There were people picking through it. There were people picking through everything.

The Guild of Alchemists’ hall was worse, and he did not understand why until a woman told him.

“Four hundred in there,” she said. She was sitting on the steps with her hands in her lap and she said it to him without appearing to particularly notice who he was. “The bald one put us in. Went down the cellar himself with a lantern to be sure of it — wouldn’t take the guildmen’s word, he said. Bricked up since Aerys, he said, and he’d seen it with his own eyes.” She looked up at the hall. The roof was gone; the dragon had taken it off on the way down and the interior was open to the sky and blackened. “Best place in the city, he said. And he was right, my lord. He was right about the cellar.”

Tyrion could not think of one word to say to her.

\textit{You gave four hundred people the safest cellar in King’s Landing,} he thought, walking on. \textit{You checked it personally. And a dead dragon fell on the roof.}

He kept walking. That was the thing about the whole of that first day: there was nowhere to stop.

He counted, because counting was what he had. The Mud Gate quarter, mostly standing, mostly emptied. Flea Bottom, sections of it simply absent. The Dragonpit intact and full of perhaps three thousand people who had climbed into the seating tiers on Davos’s instruction and had not yet been persuaded to come down. Cellars open everywhere — that was the thing that got into him worst, worse than the burnt streets. Every fourth building in the lower city had a hole in its floor with the flagstones pushed up from underneath.

\textit{Nobody is going to live in a house with a cellar again in this city for a generation. That’s not superstition. That’s a housing crisis.}

\sceneBreak

The Faith sent a septon called Cellador, who was younger than Tyrion expected and considerably harder.

“The Seven receive the dead,” he said. “In the earth, in consecrated ground, with the rites. That is not custom, my lord, and it is not preference. It is the whole of what the Faith is for. You are asking me to stand aside while thirty thousand souls are denied it.”

“I’m asking you to tell me where to put them.”

“In the ground.”

“The ground is frozen four feet down.”

“Then the Seven will be patient. The dead are patient.” Cellador’s hands were folded and he did not raise his voice at any point in the entire discussion, which Tyrion found much more difficult to argue with than shouting. “You will wait for the thaw, and you will bury them properly, and the city will endure the waiting because the alternative is that we become the sort of people who burn their mothers.”

“When it thaws, Septon, thirty thousand bodies in a city with two working wells will kill more people than the dead did.” Tyrion heard his own voice fraying and did nothing about it. “It’ll be the bloody flux by the second week and it will take the children first, and you will bury them too, in the ground, with the rites, and there will be a great many more of them than there needed to be.”

“That may be so.”

“\textit{That may be so?}”

“I did not say I was indifferent to it, my lord. I said it may be so.” Cellador looked at him steadily. “You are a clever man and you have spent this conversation trying to make me choose between the rites and the flux, and I understand why. But I would ask you to consider what the people out there have just been through, and then tell me honestly which of us is going to have more difficulty being obeyed.”

And that, Tyrion thought afterward, was the moment he understood he was not dealing with a fool, and that the problem was very much worse than the arithmetic.

Because the septon was right. The city would not bury its dead. Not out of piety — the opposite. Every soul in King’s Landing had spent two days watching their own grandmothers come up through the floor, and there was not a power on earth that was going to get them to put anybody in the ground again. They would burn. They would insist on burning.

And the Faith would forbid it. And there was no wood.

“There’s wildfire,” said Davos, from the end of the table, in the flat voice of a man who had been thinking about it for some hours and did not want to be the one who said it.

Nobody spoke for a while after that.

\sceneBreak

The last thing that happened that day was that Grey Worm came to see him.

He came alone, which Tyrion noted, and stood rather than sitting, which Tyrion also noted.

“My men are in the streets,” he said. “Eight thousand and some. They are being fed from stores that are not theirs, in a city that has no ruler, and they have been asked to do nothing for two days.” A pause. “I have kept them in order. I want you to know that this has required work.”

“I do know it.”

“There are three hundred of them in the Dragonpit at present who are preventing three thousand frightened people from coming down into streets where there is nothing to eat. Yesterday two of my men were struck with stones.” Grey Worm’s face gave up absolutely nothing. “I did not permit them to respond. Today it will be more than stones, and I will not permit them to respond then either, and I am telling you this now, before it happens, because I want it understood afterward that it was a choice and not an accident.”

“Thank you.”

“It is not a courtesy, Lord Tyrion. It is a warning.” He turned to go, and then stopped. “The queen is going to be asked what her army is for. Soon. By somebody. It would be better for everyone in this city if she has an answer before she is asked, and at present she cannot sit up.”

Then he went, and Tyrion sat alone in the Tower of the Hand in the dark and understood that he had just been handed, very politely, the single most dangerous fact in Westeros: that the only intact professional army on the continent was standing in a city with no government, taking orders from a man who was choosing, daily and without being required to, not to use it.

\sceneBreak

He got the wine at about the eleventh hour.

He had not been avoiding it, precisely. He had simply been busy, in the way a man is busy when being busy is doing something for him, and at the eleventh hour of the first day he ran out of things that had to be decided before morning and found himself in a room with a chair and a table and no further business.

Varys’s letters were still in his coat. He had been aware of them all day, the way a man is aware of a stone in his boot he has decided not to stop and remove. Eleven drafts, none of them finished, saying in eleven careful registers that Jon Snow was Aegon Targaryen and had the better claim and that the queen’s judgment ought to concern any man who intended to live under her.

He took them out and put them on the table and looked at them.

\textit{You were right,} he thought. \textit{That’s the difficulty, isn’t it. You were entirely right, and you were right in writing, and every man you addressed is alive and intact and better fed than we are, and in about three weeks the first of them is going to send someone here to ask a polite question.}

He could burn them. There was a candle. It would take a minute.

He poured the wine instead — the first of it, in a cup he had not intended to fill, in a room where nobody could see him do it — and he did not burn the letters, and he did not read them again either, and he sat for a long while in the dark with both of those decisions sitting in front of him on the same table.

Somewhere below, in the Great Hall, four hundred people were being kept alive by a fat boy from Horn Hill who had not slept in three days.

Somewhere in a storeroom off the kitchens, in the cold, his brother was waiting for a thaw that was going to kill this city.

Tyrion Lannister refilled the cup, and understood, with the particular clarity that came to him only at the worst hour of any given day, that the war had been considerably easier than this was going to be.